\chapter{Soporte de software}
\label{Soporte de software}

\section{Estrategia de mantenimiento de software}

Una vez implementada la plataforma de gestión de fletes, será necesario establecer un plan de mantenimiento que permita asegurar su correcto funcionamiento, mejorar continuamente y responder oportunamente a los cambios tecnológicos y a las necesidades de los usuarios.

Para ello se consideran cuatro tipos de mantenimientos definidos en la Ingeniería Software: correctivo, adaptativo, evolutivo y preventivo.

\noindent \textbf{Mantenimiento correctivo}

El mantenimiento correctivo tiene como objetivo solucionar errores detectados durante la operación del sistema que puedan afectar a la disponibilidad o el funcionamiento de esta.

En el contexto del software de fletes, este mantenimiento contempla acciones como:

\begin{itemize}
	\item Solución de fallos en seguimiento GPS en tiempo real
	\item Reparación de problemas en la carga o visualización de fotografías de respaldo.
	\item Solución de incidencias reportadas por clientes, transportistas o administradores
\end{itemize}

\noindent \textbf{Mantenimiento Adaptativo}

El mantenimiento adaptativo permitirá que la plataforma continúe operando correctamente frente a cambios del entorno tecnológico o normativo.

Entre las principales actividades se consideran:

\begin{itemize}
	\item Compatibilidad con nuevas versiones de Android e iOS.
	\item Actualización de servicios de mapas y geolocalización. 
	\item Integración con nuevos proveedores de pago electrónico.
	\item Adecuación a modificaciones legales relacionadas con protección de datos personales y comercio electrónico. 
\end{itemize}

\noindent \textbf{Mantenimiento Evolutivo} 

El mantenimiento evolutivo busca incorporar nuevas funcionalidades que agreguen valor al sistema a partir de la retroalimentación obtenida de los usuarios y de las necesidades del negocio.

Algunas mejoras planificadas son:

\begin{itemize}
	\item Incorporación de nuevos medios de pago.
	\item Mejoras en el historial de servicios.
	\item Mejoras en las notificaciones.
	\item Estadísticas avanzadas para transportistas y administradores.
	\item Expansión a nuevas comunas.
\end{itemize}


\noindent\textbf{Mantenimiento Preventivo}

El mantenimiento preventivo tiene como finalidad disminuir la probabilidad de fallos futuros mediante actividades periódicas de monitoreo y optimización.

Entre las tareas consideradas destacan:

\begin{itemize}
	\item Actualización de librerías y dependencias de software.
	\item Instalación periódica de parches de seguridad.
	\item Respaldos automáticos de bases de datos.
	\item Revisión de registros (logs) para detectar anomalias.
\end{itemize}


\section{Soporte técnico a los usuarios}

Con el propósito de ofrecer una atención eficiente y mantener la satisfacción de clientes y transportistas, la plataforma dispondrá de distintos canales de soporte.

\noindent\textbf{Centro de ayuda}

La aplicación incorpora una sección de preguntas (FAQ), donde los usuarios podrán resolver sus dudas relacionadas con:

\begin{itemize}
	\item Registro de cuentas.
	\item Solicitud de servicios.
	\item Metodos de pago.
	\item Seguimiento del traslado.
	\item Gestión de reclamos.
	\item Uso del sistema de fotografías.
\end{itemize}

\noindent\textbf{Sistema de tickets}

Los usuarios podrán reportar problemas mediante un formulario integrado en la aplicación. Cada reporte generará automáticamente un número de seguimiento que permitirá consultar el estado de la solicitud en cualquier momento.Además las incidencias serán calificadas según su nivel de prioridad (Baja, Media, Alta, Crítica), con el fin de optimizar su gestión. El tiempo de respuesta dependerá de la prioridad asignada, atendiendo primero las incidencias críticas por su mayor impacto en el funcionamiento del sistema, mientras que las de prioridad baja tendrán un tiempo de respuesta más amplio. 

\noindent\textbf{Chat de Soporte}

Se implementará un chat de asistencia con algún agente especializado de nuestra empresa, para resolver consultas simples relacionadas con el uso de la plataforma, permitiendo una atención rápida sin necesidad de abandonar la aplicación.

\noindent\textbf{Correo electrónico}

Para incidencias complejas, envío de documentación o apelaciones, los usuarios podrán comunicarse mediante correo electrónico con el equipo de soporte técnico.

\noindent\textbf{Gestión de reclamos}

La plataforma dispondrá de un procedimiento formal para resolver conflictos entre clientes y transportistas.

Cuando exista un reclamo, el administrador podrá revisar:

\begin{itemize}
	\item Fotografías registradas antes y después del traslado.
	\item Historial del servicio realizado
	\item Conversaciones del chat interno entre cliente y transportista
	\item Estado del servicio y registros asociados
\end{itemize}

\section{Escalabilidad y Adaptabilidad del Sistema}

\noindent\textbf{Escalabilidad}

Para soportar un aumento continuo de usuarios se consideran las siguientes estrategias:

\begin{itemize}
	\item Infraestructura basada en servicios en la nube.
	\item Balanceo de carga entre servidores.
	\item Bases de datos escalables.
	\item Almacenamiento distribuido para fotografías y archivos.
	\item Monitoreo permanente del consumo de recursos
\end{itemize}

\noindent\textbf{Adaptabilidad}

El diseño modular del sistema facilitará incorporar nuevas funcionalidades sin modificar la arquitectura principal.

Entre las futuras ampliaciones destacan:

\begin{itemize}
	\item Expansión del servicio a nuevas comunas y regiones del país.
	\item Incorporación de transporte de cargas de mayor volumen.
	\item Integración con billeteras digitales y nuevos medios de pago.
\end{itemize}

La utilización de una arquitectura modular y escalable permitirá que la plataforma evolucione conforme aumente la demanda, manteniendo elevados estándares de disponibilidad, rendimiento y calidad del servicio.
